% ============================================================
% 纯答案册·课时03-方向向量与法向量 —— 工具/答案抽册器.py 抽册生成件（勿手改；第三档＝纯答案单独成册）
% 源件（只读）：C:/提示词/工作区/M3-第2章量产0913/成卷/练习件/课时03-方向向量与法向量/main.tex（md5 fcf2aa27b35245cedbf2321523b13f34）
%% 口径：题面不重印；逐题＝题号(印面号同正册)＋[答案]＋[详解]，值/详解照源件逐字
%       （钉值门硬断言：册值≡正册件内值≡值台账值tex，逐字全等）。册序＝正册装配序＝印面号 1..16 升序（同序同号）；键序断言基准＝题面库冻结 manifest canonical 键序。
%% 三档导出：整体 main-true／纯题 main-false（在产）｜本册＝纯答案册（本件）
%% 双档：ansbook-true.tex＝详解印本；ansbook-false.tex＝纯值速查本（详解不印）
% 编译：xelatex <壳> ×2，cwd＝本目录；sty＝qp-m3.sty 本地挂载（md5 7c3930362be8a0a2bdf21bbf8ac16573，锁校验）
% ============================================================
\documentclass[fontset=none]{ctexart}
\usepackage{qp-m3}
% —— 印面逐字符号路由补钉（承源件；− 走 TNR、▱ NSC 子块；禁改 sty 保 md5） ——
\xeCJKDeclareCharClass{Default}{"2212}%
\setCJKfamilyfont{nscblk}[Path=C:/提示词/工作区/字替对照-0909/variantF/fonts/,BoldFont=NSC-w600.ttf]{NSC-w500.ttf}%
\xeCJKDeclareSubCJKBlock{nscblk}{"25B1}%
\newcommand{\pxparallelogram}{{\CJKfamily{nscblk}▱}}%
% —— 断行微调（承母版实测档，三零门承重；\emergencystretch 不另设） ——
\xeCJKsetup{CJKglue={\hskip 0.10em plus 0.08em minus 0.05em}}%
% —— 纯答案册全线括线（长详解可跨栏断，承练习件全线括线口径；灰底盒不可跨栏断） ——
\ansblockgrayfalse
% —— 双档开关（false 档＝纯值速查：答案印、详解不印；件面开关，不动 sty 本体） ——
\newif\ifansbookdetail
\ifdefined\ansbookdetail
  \ifnum\ansbookdetail=0 \ansbookdetailfalse\else\ansbookdetailtrue\fi
\else
  \ansbookdetailtrue
\fi
% —— 页脚件名（承源件章名；件名＝<件型>答案册） ——
\renewcommand{\qpzhangming}{第二章\quad 平面解析几何}%
\renewcommand{\qpjianming}{练习件答案册}%

\begin{document}
\keshi{第3课时\quad 方向向量与法向量}
\par\glueguard{1}\addvspace{2pt}\noindent\biaoqian{【纯答案册】}{\fangsong 题面不重印\quad 印面号与正册同号同序}\par\addvspace{2pt}

\begin{multicols}{2}
\emergencystretch=1em

\begin{ansblock}[2章-练-课时03-E1]
% ans:2章-练-课时03-E1
\ansitem{1}{2；4/3}
\ifansbookdetail
\ansnote{详解}{由方向向量\((2,4)\)得直线\(l\)的斜率 \(k=4/2=2\)．\par\noindent 由两点斜率公式 \(k=(m-(-2))/(3-m)=(m+2)/(3-m)\)（\(m\neq 3\)），令\((m+2)/(3-m)=2\)，得\(m+2=6-2m\)，\(3m=4\)，\(m=4/3\)．\par\noindent 故斜率为2，\(m=4/3\)．}
\fi
\end{ansblock}

\begin{ansblock}[2章-练-课时03-E2]
% ans:2章-练-课时03-E2
\ansitem{2}{(1) 1或2；(2) 4/3}
\ifansbookdetail
\ansnote{详解}{(1) \(k_{AC}=[4-(-m-3)]/(-1-m)=(m+7)/(-m-1)\)（\(m\neq -1\)），\(k_{BC}=(4-m+1)/(-1-2)=(5-m)/(-3)\)．\par\noindent 由\(k_{AC}=3k_{BC}\)：\((m+7)/(-m-1)=3\cdot(5-m)/(-3)=-(5-m)=m-5\)，\par\noindent 去分母得 \(m+7=(m-5)(-m-1)=-m^{2}+4m+5\)，即 \(m^{2}-3m+2=0\)，解得\(m=1\)或\(m=2\)，经检验均使两斜率存在，符合题意．\par\noindent (2) 设\(l_1\)的倾斜角为\(\alpha\)，则\(l_2\)的倾斜角为\(2\alpha\)．由方向向量\(\overrightarrow{n}=(2,1)\)得 \(\tan\alpha=1/2\)，\par\noindent 故\(l_2\)的斜率 \(\tan 2\alpha=2\tan\alpha/(1-\tan^{2}\alpha)=1/(1-1/4)=4/3\)．}
\fi
\end{ansblock}

\begin{ansblock}[2章-练-课时03-E3]
% ans:2章-练-课时03-E3
\ansitem{3}{方向向量如(1,−1)（答案不唯一）；k＝−1；θ＝135°}
\ifansbookdetail
\ansnote{详解}{\(\overrightarrow{AB}=(-5+2,\ 3-0)=(-3,3)\)是直线\(l\)的一个方向向量（与之共线的非零向量如\((1,-1)\)也可）．\par\noindent 斜率 \(k=(3-0)/(-5-(-2))=3/(-3)=-1\)（也可由方向向量\((1,-1)\)得\(k=-1/1=-1\)）．\par\noindent 由\(\tan\theta=-1\)且\(\theta\in[0^\circ,180^\circ)\)，得\(\theta=135^\circ\)．}
\fi
\end{ansblock}

\begin{ansblock}[2章-练-课时03-E4]
% ans:2章-练-课时03-E4
\ansitem{4}{A、B、C不共线；A、B、D共线．}
\ifansbookdetail
\ansnote{详解}{\(\overrightarrow{AB}=(0+1,-1+3)=(1,2)\)，\(\overrightarrow{AC}=(1+1,2+3)=(2,5)\)．\par\noindent 因为\(1\times 5-2\times 2=1\neq 0\)，\(\overrightarrow{AB}\)与\(\overrightarrow{AC}\)不共线，且两向量有公共起点\(A\)，所以\(A\)、\(B\)、\(C\)不共线．\par\noindent \(\overrightarrow{AD}=(3+1,5+3)=(4,8)=4\overrightarrow{AB}\)，\(\overrightarrow{AB}\)与\(\overrightarrow{AD}\)共线且有公共起点\(A\)，所以\(A\)、\(B\)、\(D\)共线．}
\fi
\end{ansblock}

\begin{ansblock}[2章-练-课时03-E5]
% ans:2章-练-课时03-E5
\ansitem{5}{真命题．证明见详解．}
\ifansbookdetail
\ansnote{详解}{必要性：若\(A\)、\(B\)、\(C\)共线，则直线上的向量\(\overrightarrow{AB}\)、\(\overrightarrow{BC}\)及与直线平行的向量都是该直线的方向向量，故\(\overrightarrow{AB}\)与\(\overrightarrow{BC}\)共线．\par\noindent 充分性：若\(\overrightarrow{AB}\)与\(\overrightarrow{BC}\)共线，则直线\(AB\)与直线\(BC\)的方向相同或相反；又两直线过公共点\(B\)，由过一点且方向确定的直线唯一，直线\(AB\)与直线\(BC\)是同一条直线，故\(A\)、\(B\)、\(C\)三点共线．\par\noindent 综上命题为真．}
\fi
\end{ansblock}

\begin{ansblock}[2章-练-课时03-E6]
% ans:2章-练-课时03-E6
\ansitem{6}{(1) y＝−2x−4（即2x＋y＋4＝0）；(2) 当u≠0时为y−y₀＝(v/u)(x−x₀)；当u＝0时为x＝x₀．}
\ifansbookdetail
\ansnote{详解}{(1) 方向向量\(a=(1,-2)\)的横坐标\(1\neq 0\)，故斜率\(k=-2/1=-2\)，由点斜式得 \(y-(-2)=-2[x-(-1)]\)，即 \(y=-2x-4\)．\par\noindent (2) 设\(Q(x,y)\)是直线\(l\)上任意一点，则\(\overrightarrow{PQ}=(x-x_0,\ y-y_0)\)与方向向量\(a=(u,v)\)共线，得 \(v(x-x_0)-u(y-y_0)=0\)．\par\noindent 当\(u\neq 0\)时化为 \(y-y_0=(v/u)(x-x_0)\)；当\(u=0\)时（此时\(v\neq 0\)）得 \(x-x_0=0\)，即\(x=x_0\)（斜率不存在的竖直线）．}
\fi
\end{ansblock}

\begin{ansblock}[2章-练-课时03-E7]
% ans:2章-练-课时03-E7
\ansitem{7}{(1) y＝3x；(2) u(x−x₀)＋v(y−y₀)＝0，即ux＋vy−ux₀−vy₀＝0．}
\ifansbookdetail
\ansnote{详解}{(1) 设\(Q(x,y)\)是直线\(l\)上任意一点，则\(\overrightarrow{PQ}=(x-1,\ y-3)\)与法向量\(a=(-3,1)\)垂直，\(a\cdot\overrightarrow{PQ}=0\)，\par\noindent 即 \(-3(x-1)+1\cdot(y-3)=0\)，化简得 \(y=3x\)．\par\noindent (2) 同法：\(a\cdot\overrightarrow{PQ}=u(x-x_0)+v(y-y_0)=0\)（\(u\)、\(v\)不全为0，因法向量是非零向量）．这就是直线的点法式方程；\par\noindent 当\(v\neq 0\)时可化为斜截形 \(y=-(u/v)x+(ux_0+vy_0)/v\)，斜率\(k=-u/v\)（与“方向向量\((v,-u)\)读斜率\(-u/v\)”一致）．}
\fi
\end{ansblock}

\begin{ansblock}[2章-练-课时03-E8]
% ans:2章-练-课时03-E8
\ansitem{8}{一定．理由见详解．}
\ifansbookdetail
\ansnote{详解}{设\((x_0,y_0)\)是\(l\)的一个方向向量．由方向向量非零知\((x_0,y_0)\neq(0,0)\)，从而\((-y_0,x_0)\neq(0,0)\)．\par\noindent 又 \((x_0,y_0)\cdot(-y_0,x_0)=-x_0y_0+y_0x_0=0\)，即\((-y_0,x_0)\)与\(l\)的方向向量垂直，于是以\((-y_0,x_0)\)为方向向量的有向线段所在直线与\(l\)垂直，\par\noindent 由法向量的定义，\((-y_0,x_0)\)是\(l\)的一个法向量．\par\noindent 反之，若\((-y_0,x_0)\)为\(l\)的方向向量，同理\((0,0)\neq(-y_0,x_0)\)\(\Longleftrightarrow\)\((x_0,y_0)\neq(0,0)\)，且两向量数量积为0，故\((x_0,y_0)\)一定是\(l\)的法向量．\par\noindent 所以结论成立（教材正文中“特别地，当\(x_0\)与\(y_0\)不全为0时…”的前提恰好被“方向向量非零”自动满足）．}
\fi
\end{ansblock}

\begin{ansblock}[2章-练-课时03-E9]
% ans:2章-练-课时03-E9
\ansitem{9}{(1) y＝2x−11；(2) y＝−(4/3)x＋5（即4x＋3y−15＝0）}
\ifansbookdetail
\ansnote{详解}{(1) \(n=(1,2)\)横坐标不为0，\(k=2/1=2\)，由点斜式 \(y+5=2(x-3)\)，即 \(y=2x-11\)．\par\noindent (2) \(n=(3,-4)\)横坐标不为0，\(k=-4/3\)，由点斜式 \(y-5=-(4/3)(x-0)\)，即 \(y=-(4/3)x+5\)．}
\fi
\end{ansblock}

\begin{ansblock}[2章-练-课时03-E10]
% ans:2章-练-课时03-E10
\ansitem{10}{(1) 3x−4y＋5＝0；(2) 3x＋4y−5＝0}
\ifansbookdetail
\ansnote{详解}{(1) 设\(Q(x,y)\)在\(l\)上，\(v\cdot\overrightarrow{PQ}=3(x-1)-4(y-2)=0\)，化简得 \(3x-4y+5=0\)．\par\noindent (2) 同法：\(3(x+1)+4(y-2)=0\)，化简得 \(3x+4y-5=0\)．}
\fi
\end{ansblock}

\begin{ansblock}[2章-练-课时03-E11]
% ans:2章-练-课时03-E11
\ansitem{11}{证明见详解．}
\ifansbookdetail
\ansnote{详解}{\(\overrightarrow{AB}=(3-1,3-4)=(2,-1)\)，故直线\(AB\)的一个方向向量为\(a=(2,-1)\)；\(\overrightarrow{BC}=(4-3,5-3)=(1,2)\)．\par\noindent 取\(v=(1,2)\)，验证\(v\)与\(a\)垂直：\(v\cdot a=1\times 2+2\times(-1)=0\)，且\(v\neq 0\)，所以\(v\)是直线\(AB\)的一个法向量．\par\noindent 而\(\overrightarrow{BC}=v\)，即直线\(BC\)的方向向量就是直线\(AB\)的法向量，由法向量的定义，直线\(BC\)与直线\(AB\)垂直，故\(AB\perp BC\)．}
\fi
\end{ansblock}

\begin{ansblock}[2章-练-课时03-E12]
% ans:2章-练-课时03-E12
\ansitem{12}{B}
\ifansbookdetail
\ansnote{详解}{\(AB=(3-1,\ 6-2)=(2,4)=2(1,2)\)，故 \((1,2)\) 为方向向量，选 B。}
\fi
\end{ansblock}

\begin{ansblock}[2章-练-课时03-E13]
% ans:2章-练-课时03-E13
\ansitem{13}{2x−y−5=0}
\ifansbookdetail
\ansnote{详解}{法向量 \((2,-1)\)，故 \(l\) 具形式 \(2x-y+C=0\)。代 \(P(3,1)\)：\(6-1+C=0\)，即 \(C=-5\)。故 \(2x-y-5=0\)。}
\fi
\end{ansblock}

\begin{ansblock}[2章-练-课时03-E14]
% ans:2章-练-课时03-E14
\ansitem{14}{(1) n₁=(3,−1)（任一非零倍数均可）　(2) a=2/3}
\ifansbookdetail
\ansnote{详解}{(1) 由一般式系数直接读出 \(n_1=(3,-1)\)。(2) \(n_2=(a,2)\)。同一平面内每条直线的法向量与其方向向量互相垂直，将两直线绕任一点同旋 \(90^\circ\)，方向向量恰转到各自法向量方向且夹角不变，\par\noindent 故 \(l_1\perp l_2\)\(\Longleftrightarrow\)\(n_1\perp n_2\)\(\Longleftrightarrow\)\(n_1\cdot n_2=0\)\(\Longleftrightarrow\)\(3a+(-1)\times 2=0\)\(\Longleftrightarrow\)\(a=2/3\)。}
\fi
\end{ansblock}

\begin{ansblock}[2章-练-课时03-E15]
% ans:2章-练-课时03-E15
\ansitem{15}{m=0}
\ifansbookdetail
\ansnote{详解}{直线 \(Ax+By+C=0\) 的方向向量可取 \((B,\ -A)\)，故 \(l\) 的方向向量 \(t=(2,\ -(m+1))\)。\par\noindent \(t\parallel v\)（横坐标已同为 2），即 \(-(m+1)=m-1\)，\par\noindent 解得 \(2m=0\)，故 \(m=0\)。}
\fi
\end{ansblock}

\begin{ansblock}[2章-练-课时03-E16]
% ans:2章-练-课时03-E16
\ansitem{16}{−√3/3}
\ifansbookdetail
\ansnote{详解}{斜率 \(k=\sqrt{3}/1=\sqrt{3}\)，\(l\): \(y=\sqrt{3}x+1\)。令 \(y=0\)：\(x=-1/\sqrt{3}=-\sqrt{3}/3\)。}
\fi
\end{ansblock}

% ---- 尾块（\tailfill[30mm] 定高参——钉档 §三.3：无参在平衡栏呈「内容尾随框」，贴栏底须定高参；实测无参致末页平衡 Overfull\vbox 12.99pt，定高 30mm 三零） ----
\tailfill[30mm]

\end{multicols}
\end{document}
